\documentclass[10pt,a4paper]{amsart}
\usepackage{amsmath,amssymb,amsthm,graphicx}
% For commutative diagrams we use Xy-pic
\usepackage[all]{xy}

% We use 2cell for 2-commutative diagrams.
\xyoption{2cell}
% Theorem environments.
%
\theoremstyle{plain}
\newtheorem{theorem}[subsection]{Theorem}
\newtheorem{proposition}[subsection]{Proposition}
\newtheorem{lemma}[subsection]{Lemma}

\theoremstyle{definition}
\newtheorem{definition}[subsection]{Definition}
\newtheorem{example}[subsection]{Example}
\newtheorem{exercise}[subsection]{Exercise}
\newtheorem{situation}[subsection]{Situation}

\theoremstyle{remark}
\newtheorem{remark}[subsection]{Remark}
\newtheorem{remarks}[subsection]{Remarks}

\numberwithin{equation}{subsection}

% Macros
%
\def\lim{\mathop{\mathrm{lim}}\nolimits}
\def\colim{\mathop{\mathrm{colim}}\nolimits}
\def\Spec{\mathop{\mathrm{Spec}}}
\def\Hom{\mathop{\mathrm{Hom}}\nolimits}
\def\Ext{\mathop{\mathrm{Ext}}\nolimits}
\def\SheafHom{\mathop{\mathcal{H}\!\mathit{om}}\nolimits}
\def\SheafExt{\mathop{\mathcal{E}\!\mathit{xt}}\nolimits}
\def\Sch{\mathit{Sch}}
\def\Mor{\mathop{\mathrm{Mor}}\nolimits}
\def\Ob{\mathop{\mathrm{Ob}}\nolimits}
\def\Sh{\mathop{\mathit{Sh}}\nolimits}
\def\NL{\mathop{N\!L}\nolimits}
\def\CH{\mathop{\mathrm{CH}}\nolimits}
\def\proetale{{pro\text{-}\acute{e}tale}}
\def\etale{{\acute{e}tale}}
\def\QCoh{\mathit{QCoh}}
\def\Ker{\mathop{\mathrm{Ker}}}
\def\Im{\mathop{\mathrm{Im}}}
\def\Coker{\mathop{\mathrm{Coker}}}
\def\Coim{\mathop{\mathrm{Coim}}}

% Boxtimes
%
\DeclareMathSymbol{\boxtimes}{\mathbin}{AMSa}{"02}

%
% Macros for moduli stacks/spaces
%
\def\QCohstack{\mathcal{QC}\!\mathit{oh}}
\def\Cohstack{\mathcal{C}\!\mathit{oh}}
\def\Spacesstack{\mathcal{S}\!\mathit{paces}}
\def\Quotfunctor{\mathrm{Quot}}
\def\Hilbfunctor{\mathrm{Hilb}}
\def\Curvesstack{\mathcal{C}\!\mathit{urves}}
\def\Polarizedstack{\mathcal{P}\!\mathit{olarized}}
\def\Complexesstack{\mathcal{C}\!\mathit{omplexes}}
% \Pic is the operator that assigns to X its picard group, usage \Pic(X)
% \Picardstack_{X/B} denotes the Picard stack of X over B
% \Picardfunctor_{X/B} denotes the Picard functor of X over B
\def\Pic{\mathop{\mathrm{Pic}}\nolimits}
\def\Picardstack{\mathcal{P}\!\mathit{ic}}
\def\Picardfunctor{\mathrm{Pic}}
\def\Deformationcategory{\mathcal{D}\!\mathit{ef}}
\usepackage{cleveref}
\crefname{originalequation}{original equation}{original equations}
\Crefname{originalequation}{Original equation}{Original equations}
\makeatletter
\newlabel{equation-tame-symbol}{{0EAQ}{source}{}{sourcepointer-1}{}}
\newlabel{equation-tame-symbol@cref}{{[originalequation][0][]0EAQ}{[1][0][]source}{}{}{}}
\makeatother
\crefname{originallemma}{original lemma}{original lemmas}
\Crefname{originallemma}{Original lemma}{Original lemmas}
\makeatletter
\newlabel{lemma-prepare-tame-symbol}{{0EAG}{source}{}{sourcepointer-2}{}}
\newlabel{lemma-prepare-tame-symbol@cref}{{[originallemma][0][]0EAG}{[1][0][]source}{}{}{}}
\makeatother
\crefname{originallemma}{original lemma}{original lemmas}
\Crefname{originallemma}{Original lemma}{Original lemmas}
\makeatletter
\newlabel{fields-lemma-trace-and-norm-tower}{{0BIJ}{source}{}{sourcepointer-3}{}}
\newlabel{fields-lemma-trace-and-norm-tower@cref}{{[originallemma][0][]0BIJ}{[1][0][]source}{}{}{}}
\makeatother
\crefname{originallemma}{original lemma}{original lemmas}
\Crefname{originallemma}{Original lemma}{Original lemmas}
\makeatletter
\newlabel{lemma-compare-periodic-lengths}{{0EA9}{source}{}{sourcepointer-4}{}}
\newlabel{lemma-compare-periodic-lengths@cref}{{[originallemma][0][]0EA9}{[1][0][]source}{}{}{}}
\makeatother
\crefname{originalequation}{original equation}{original equations}
\Crefname{originalequation}{Original equation}{Original equations}
\makeatletter
\newlabel{equation-multiplicity-coker-ker}{{0EA6}{source}{}{sourcepointer-5}{}}
\newlabel{equation-multiplicity-coker-ker@cref}{{[originalequation][0][]0EA6}{[1][0][]source}{}{}{}}
\makeatother
\crefname{originallemma}{original lemma}{original lemmas}
\Crefname{originallemma}{Original lemma}{Original lemmas}
\makeatletter
\newlabel{algebra-lemma-pushdown-module}{{02M0}{source}{}{sourcepointer-6}{}}
\newlabel{algebra-lemma-pushdown-module@cref}{{[originallemma][0][]02M0}{[1][0][]source}{}{}{}}
\makeatother
\title{The tame symbol is well defined independently of the factorizations and of the chosen extension}
\author{The Stacks Project Authors}
\date{stacks-project:chow.tex:local-snapshot-20261012; GNU FDL 1.2 or later}
\begin{document}
\maketitle
\noindent\textit{Editorial note.} All delivered mathematical text is the author's own, copied from the pinned Stacks Project chapter file chow.tex; no statement or proof step was rewritten, reordered or removed. The author's notation (1)-(6) list inside the proof and the phrase 'This proves the desired equality up to signs' are reproduced verbatim.\par\medskip
\noindent\textit{Editorial note.} source\_id and source\_version carry the current authoritative snapshot label of this source, 20261012, read from sources/stacks-chow/license-observation.json and SOURCE.json at the moment of opening.\par\medskip
\noindent\textit{Editorial note.} The preamble is the chapter's own native preamble excerpt used for this source: lines 14-18 (Xy-pic setup) and 154-217 (theorem declarations and mathematical macros) of sources/stacks-chow/preamble.tex. The quoted class line and the xr-hyper externaldocument block are not included because the assembler supplies a plain amsart wrapper; no macro was invented or redefined.\par\medskip
\noindent\textit{Editorial note.} Three of the referenced results are same-chapter and are located at their \textbackslash{}label lines inside the pinned chapter file; the two cross-chapter pointers are located in the project's own tag list staged in the snapshot (sources/stacks-chow/tags.txt).\par\medskip
\noindent\textit{Allowed standard prerequisites (editorial).}\begin{itemize}
\item The displayed formula defining the tame symbol, and the preparation lemma producing local factorizations. Same chapter, the displayed formula labelled equation-tame-symbol (source line 820, Stacks tag 0EAQ) and Lemma labelled lemma-prepare-tame-symbol (source line 642, Stacks tag 0EAG); named in the counted statement at line 831 and invoked at proof line 865. Supply the formula whose well-definedness is asserted and the class of extensions B over which the reader is allowed to compute; the two independence arguments are the counted content.
\item Transitivity of the norm in a tower of field extensions. Fields, Lemma labelled fields-lemma-trace-and-norm-tower (canonical Stacks tag 0BIJ); invoked at proof lines 875-876. Supplies the reduction of the independence of B to the numerical identity lambda\textasciicircum{}m = prod Norm(lambda)\textasciicircum{}\{m\_k\}.
\item Length bookkeeping for a finite extension of one-dimensional local rings. Same chapter, Lemma lemma-compare-periodic-lengths (source line 283, Stacks tag 0EA9) with the equation equation-multiplicity-coker-ker (source line 203, Stacks tag 0EA6), and Algebra, Lemma algebra-lemma-pushdown-module (canonical Stacks tag 02M0); invoked at proof lines 894-898. Supply m = length\_B(C/pi C) = sum m\_k [kappa\_k : kappa], the equality that makes the displayed norm identity true.
\end{itemize}\medskip
\ifdefined\hypertarget\hypertarget{sourcepointer-1}{}\fi
\noindent\textit{Source reference.} Original equation 0EAQ: sources/stacks-chow/pinned-original-source.tex, source lines 820--820. Same-chapter pointer to the displayed definition of the tame symbol, named in the counted statement at line 831; both independence proofs are retained in the counted span.\par\medskip
\ifdefined\hypertarget\hypertarget{sourcepointer-2}{}\fi
\noindent\textit{Source reference.} Original lemma 0EAG: sources/stacks-chow/pinned-original-source.tex, source lines 642--642. Same-chapter pointer invoked at proof line 865 for the class of admissible extensions; the reduction to the norm identity is retained in the counted span.\par\medskip
\ifdefined\hypertarget\hypertarget{sourcepointer-3}{}\fi
\noindent\textit{Source reference.} Original lemma 0BIJ: sources/stacks-chow/tags.txt, source lines 14492--14492. Cross-chapter standard prerequisite invoked once at proof lines 875-876 for transitivity of norms.\par\medskip
\ifdefined\hypertarget\hypertarget{sourcepointer-4}{}\fi
\noindent\textit{Source reference.} Original lemma 0EA9: sources/stacks-chow/pinned-original-source.tex, source lines 283--283. Same-chapter pointer invoked at proof line 895 for the length comparison m = length\_B(C/pi C); the final assembly of the identity is retained in the counted span.\par\medskip
\ifdefined\hypertarget\hypertarget{sourcepointer-5}{}\fi
\noindent\textit{Source reference.} Original equation 0EA6: sources/stacks-chow/pinned-original-source.tex, source lines 203--203. Same-chapter pointer to the defining display of e\_R, invoked together with the comparison lemma at proof line 896.\par\medskip
\ifdefined\hypertarget\hypertarget{sourcepointer-6}{}\fi
\noindent\textit{Source reference.} Original lemma 02M0: sources/stacks-chow/tags.txt, source lines 3427--3427. Cross-chapter standard prerequisite invoked once at proof line 898 for length\_B(C/pi C) = sum m\_k [kappa\_k : kappa], the last numerical input.\par\medskip
\begin{lemma}
\label{lemma-well-defined-tame-symbol}
The formula (\ref{equation-tame-symbol}) determines a
well defined element of $\kappa(\mathfrak m)^*$. In other words, the
right hand side does not depend on the choice of the
local factorizations or the choice of $B$.
\end{lemma}

\begin{proof}
Independence of choice of factorizations. Suppose we have
a Noetherian $1$-dimensional local ring $B$, elements $a_1, a_2 \in B$,
and nonzerodivisors $\pi, \theta$ such that we can write
$$
a_1 = u_1 \pi^{e_1} = v_1 \theta^{f_1},\quad
a_2 = u_2 \pi^{e_2} = v_2 \theta^{f_2}
$$
with $e_i, f_i \geq 0$ integers and $u_i, v_i$ units in $B$.
Observe that this implies
$$
a_1^{e_2} = u_1^{e_2}u_2^{-e_1}a_2^{e_1},\quad
a_1^{f_2} = v_1^{f_2}v_2^{-f_1}a_2^{f_1}
$$
On the other hand, setting
$m = \text{length}_B(B/\pi B)$ and $k = \text{length}_B(B/\theta B)$
we find $e_2 m = \text{length}_B(B/a_2 B) = f_2 k$.
Expanding $a_1^{e_2m} = a_1^{f_2 k}$ using the above we find
$$
(u_1^{e_2}u_2^{-e_1})^m =  (v_1^{f_2}v_2^{-f_1})^k
$$
This proves the desired equality up to signs. To see the signs
work out we have to show $me_1e_2$ is even if and only if
$kf_1f_2$ is even. This follows as both $me_2 = kf_2$ and
$me_1 = kf_1$ (same argument as above).

\medskip\noindent
Independence of choice of $B$. Suppose given two extensions
$A \subset B$ and $A \subset B'$ as in Lemma \ref{lemma-prepare-tame-symbol}.
Then
$$
C = (B \otimes_A B')/(\mathfrak m\text{-power torsion})
$$
will be a third one. Thus we may assume we have
$A \subset B \subset C$ and factorizations over the
local rings of $B$ and we have to show that using
the same factorizations over the local rings of $C$
gives the same element of $\kappa(\mathfrak m)$.
By transitivity of norms
(Fields, Lemma \ref{fields-lemma-trace-and-norm-tower})
this comes down to the following problem:
if $B$ is a Noetherian local ring of dimension $1$
and $\pi \in B$ is a nonzerodivisor, then
$$
\lambda^m = \prod \text{Norm}_{\kappa_k/\kappa}(\lambda)^{m_k}
$$
Here we have used the following notation:
(1) $\kappa$ is the residue field of $B$,
(2) $\lambda$ is an element of $\kappa$,
(3) $\mathfrak m_k \subset C$ are the maximal ideals of $C$,
(4) $\kappa_k = \kappa(\mathfrak m_k)$ is the residue field of
$C_k = C_{\mathfrak m_k}$,
(5) $m = \text{length}_B(B/\pi B)$, and
(6) $m_k = \text{length}_{C_k}(C_k/\pi C_k)$.
The displayed equality holds because
$\text{Norm}_{\kappa_k/\kappa}(\lambda) = \lambda^{[\kappa_k : \kappa]}$
as $\lambda \in \kappa$ and because $m = \sum m_k[\kappa_k:\kappa]$.
First, we have $m = \text{length}_B(B/xB) = \text{length}_B(C/\pi C)$
by Lemma \ref{lemma-compare-periodic-lengths} 
and (\ref{equation-multiplicity-coker-ker}).
Finally, we have $\text{length}_B(C/\pi C) = \sum m_k[\kappa_k:\kappa]$
by Algebra, Lemma \ref{algebra-lemma-pushdown-module}.
\end{proof}

\end{document}
